\chapter{绪论}\label{chap:intro}

\section{二级标题}

\subsection{三级标题}

单文献引用~\cite{thrun2005probabilistic}。多文献引用~\cite{khatib1986real, fox1997dynamic}

\section{图片插入示例}

图片文件放在 \texttt{figure/} 目录下，正文中用 \verb|\includegraphics| 插入：

\begin{figure}[htb]
    \centering
    \includegraphics[width=0.85\textwidth]{example/fig_framework}
    \caption{图片题注示例}
    \label{fig:example}
\end{figure}

图引用示例：如图~\ref{fig:example}~所示。单文献引用~\cite{quigley2009ros}。

\section{表格插入示例}

\begin{table}[htb]
    \centering
    \caption{表格题注示例}
    \label{tab:example}
    \small
    \begin{tabular}{ccc}
        \toprule
        \textbf{列一} & \textbf{列二} & \textbf{列三} \\
        \midrule
        示例数据 & 示例数据 & 示例数据 \\
        示例数据 & 示例数据 & 示例数据 \\
        \bottomrule
    \end{tabular}
\end{table}

表引用示例：如表~\ref{tab:example}~所示。

\section{公式插入示例}

\begin{equation}
    \dot{\bm x} = \bm A \bm x + \bm B \bm u
    \label{eq:example}
\end{equation}

公式引用示例：如式~\eqref{eq:example}~所示。单文献引用~\cite{sutton2018reinforcement}。

公式中的矢量、矩阵等需要加粗时，请用 \verb|\bm{}|（如 \verb|\bm x|、\verb|\bm A|），
**不要使用 \verb|\boldsymbol{}|**。模板已在 \texttt{config/packages.tex} 中将
\verb|\bm| 重定义为 \verb|\symbfit|（基于 \texttt{unicode-math}），可与正文数学字体保持一致。

\section{物理单位书写示例}

正文中的数值与物理单位请统一用 \texttt{siunitx} 宏包排版，**不要手写单位字符串**。
手写单位容易出现正/斜体不分、缺少间隔、与数学变量混淆等问题（例如把毫秒写成斜体的
\textit{m s} 或在正文中直接写 ``ms''、``rad/s''、``Nm''）。

\textbf{不带数值的单位}用 \verb|\si{}| 排版，单位内部的词与词之间不要手敲，交给宏包处理：

\begin{itemize}
    \item \verb|\si{\milli\second}|：采样时间间隔等毫秒级单位
    \item \verb|\si{\radian\per\second}|：角速度单位（不要写成 ``rad/s''）
    \item \verb|\si{\newton\meter}|：力矩单位（不要写成 ``Nm''）
    \item \verb|\si{\degree}|：角度（不要写成 ``deg''）
    \item \verb|\si{\degreeCelsius}|：摄氏温度
\end{itemize}

\textbf{带数值}时，数值用正文模式书写，单位用 \verb|\si{}|，两者以 \verb|~| 相连：

\begin{itemize}
    \item \verb|10~\si{\milli\second}|：正文渲染为 10~\si{\milli\second}
    \item \verb|1.5~\si{\radian\per\second}|：正文渲染为 1.5~\si{\radian\per\second}
    \item \verb|30~\si{\degree}|：正文渲染为 30~\si{\degree}
\end{itemize}

若想由宏包一并处理“数值 + 单位”的间距，也可使用 \verb|\SI{10}{\milli\second}|。

其余约定：正文中的纯数字、百分号等使用正文模式，只有数学变量、公式与推导才进入数学模式；
科学计数法用文本形式书写，如 \verb|7.7 \texttimes\ 10\textsuperscript{-3}|。完整规范见 \texttt{docs/latex\_rules.md}。

\section{综合引用示例}

以上各节分别给出了图片（图~\ref{fig:example}~）、表格（表~\ref{tab:example}~）、公式（式~\eqref{eq:example}~）与文献引用（如~\cite{sutton2018reinforcement, quigley2009ros}~）的写法示例。

\section{修订标记示例}

审阅修改期间，可用 \verb|\rev{}|（中文/混排）与 \verb|\revEN{}|（纯英文）标出修订处，再由总开关一键显隐：

\begin{itemize}
    \item \verb|\rev{}|：开启开关时内容以黄底高亮标出，支持跨行与 \verb|\ref| 等宏，但不可含 \verb|\cite|。
          例如：\rev{这里是被标记的中文修订文字}~\cite{sutton2018reinforcement}。
    \item \verb|\revEN{}|：开启开关时英文以黄底高亮标出，内容请勿含 \verb|\cite|、\verb|\ref| 等宏。
          例如：\revEN{This is a marked English revision sentence.}
\end{itemize}

开关位于 \texttt{zjuthesis.tex} 中：\verb|\showrevtrue| 显示上述标记；
\verb|\showrevfalse| 隐藏标记并直接输出原文，可用于生成正式提交版。

